% --- LaTeX Essay Template - S. Venkatraman ---

% --- Set document class and font size ---

\documentclass[letterpaper, 11pt]{article}

% --- Package imports ---

% lipsum is just for generating placeholder text and can be removed
\usepackage{hyperref, lipsum} 

% --- Fonts (Set to Palatino or Times New Roman if desired) ---

% \usepackage{newpxtext, newpxmath}
% \usepackage{times}

% --- Page layout settiings ---

% Set page margins
\usepackage[left=1.2in, right=1.2in, top=.9in, bottom=.9in]{geometry}

% Anchor footnotes to the bottom of the page
\usepackage[bottom]{footmisc}
\usepackage{setspace}
\setstretch{1.15}
% Set line spacing
\renewcommand{\baselinestretch}{1.2}

% Use this for a one-off '*' footnote (e.g., a URL in a heading)
\newcommand{\starfootnote}[1]{%
\begingroup
\renewcommand{\thefootnote}{\fnsymbol{footnote}}%
\footnote{#1}%
\endgroup
\addtocounter{footnote}{-1}%
}

% Set spacing between paragraphs
\setlength{\parskip}{2mm}

% --- Page formatting settings ---



% --- Essay title and author ---
\title{\vspace{-1.5cm} Summary of \textit{What's Wrong with Social Science and How to Fix It: Reflections After Reading 2578 Papers}\footnote{\url{https://www.fantasticanachronism.com/p/whats-wrong-with-social-science-and-how-to-fix-it}}}

\date{\today}

% --- Main text ---
% --- Document starts here ---

\begin{document}
% Set link colors for labeled items and URLs (blue) 
\hypersetup{colorlinks=true, linkcolor=blue, urlcolor=blue}
\maketitle

The article argues that many problems in social science.
It critiques practices such as p-hacking, publication bias, and underpowered studies, which inflate false discoveries.
Even top journals publish many non-reproducible findings, and replication attempts often fail.

Many studies could rely on small samples, flexible analytic decisions, and significance thresholds that encourage false positives and lack conceptual clarity.
Publication bias favors novel, statistically significant findings, while null results and replications receive less attention.
As a result, reported effects often fail to reproduce when independently tested.
In fact, what we need is replication, robust theoratical frameworks, and transparency.
The article calls for reforms in the research system to improve the reliability and credibility of social science findings.

\end{document}
